% Append-only section fragment. Paste below the preceding section in the master document.
% This fragment intentionally contains no preamble and no document boundaries.

\section{Persona-Private Communication}
\label{sec:pilot-communication}

Pilot Fabric 0.51 converts communication from a generic agent route into a
persona-private, receipt-verified capability.  The shared television may begin a
request, but it never selects a recipient, exposes an address, or sends content.
Recipient resolution, content review, approval, and account execution occur on a
trusted private phone bound to the active identity.

\subsection{Communication State Machine}

The governed sequence is
\begin{enumerate}
  \item resolve exactly one contact from the active persona's private contact book;
  \item create an AION Native draft without requiring a paid model;
  \item display the exact recipient, channel, subject, and body on the private phone;
  \item bind approval to a hash of the recipient and content;
  \item invoke one authorized, persona-specific provider adapter;
  \item accept success only when the provider returns a verifiable message identifier;
  \item retain a normalized receipt without the raw provider response or credentials.
\end{enumerate}
Changing either the recipient or content invalidates the approval.  An adapter
exception is recorded as an unknown delivery effect requiring reconciliation;
Pilot does not retry automatically and risk duplicate delivery.

\subsection{Inbox Intelligence and Follow-Through}

Authorized incoming-message adapters may submit a bounded record containing the
provider message identifier, sender, receive time, source, and content.  Pilot
stores a deterministic local summary, source provenance, freshness at ingestion,
and a body hash; the raw body is not retained by this layer.  A private user may
convert a message into a task or a calendar proposal.  Conversion does not inherit
authority from the message and cannot change an external service silently.

After a provider-verified send, the owner may schedule a conditional follow-up
such as ``remind me tomorrow if no verified reply has arrived.''  The condition
is linked to the verified sent-message record rather than inferred from speech.

\subsection{Children, Messaging, and Calls}

External communication from a guardian-controlled child identity requires two
decisions: the configured guardian first approves the draft, then the child
identity approves the exact content from its trusted surface.  Ordinary Pilot
messages remain distinct from external email and messaging.

WhatsApp is represented as an authorized handoff boundary.  Pilot does not scrape
WhatsApp contacts, automate an unsupported consumer session, or claim delivery
without an official provider result.  A call action prepares the private contact
and route, but placing the call requires an explicit tap on an authorized phone
surface.

\subsection{Provider and OAuth Boundary}

The Google account flow uses an installed-application OAuth design with PKCE,
a one-time state value, a system browser, an exact scope review, and a ten-minute
authorization lifetime.  Fabric retains only an opaque vault reference.  Access
and refresh tokens must be stored by an authorized operating-system or production
secret vault and are never placed in the television, nodes, audit log, or Fabric
state.  Gmail delivery uses the official \texttt{users.messages.send} endpoint
and requires a returned provider message identifier.

Public distribution still requires a registered Google OAuth client, enabled
APIs, an owned production domain, privacy policy, consent-screen configuration,
and any Google verification required for the selected user-data scopes.  Those
provider approvals cannot be completed by source code alone.

\section{Calendar and Planning}
\label{sec:pilot-calendar-planning}

The calendar layer separates create, reschedule, cancel, availability, and travel
planning into distinct governed operations.  Every modification carries an exact
scope containing its title, interval, location, reminder, travel allowance, and
provider event identifier where applicable.  Private approval signs a canonical
hash of that scope.

\subsection{Conflict and Availability Privacy}

Before approval, Pilot compares the proposed interval---including an optional
travel allowance---with normalized private busy intervals.  A conflict response
contains only the conflicting start and end boundaries.  It does not disclose
the title, attendees, description, or location of the private event.  An owner
must separately confirm that an action may proceed despite a detected conflict.

Availability queries return free intervals of a selected duration and explicitly
record that private event details were not disclosed.  This contract aligns with
the provider free/busy abstraction and allows household coordination without
exposing one person's calendar contents.

\subsection{Execution and Receipts}

Create uses an exact provider insert operation.  Reschedule requires a previously
identified provider event and sends only the reviewed time change.  Cancel also
requires an exact event identifier.  Each action has a fresh approval and an
idempotent execution state.  A normalized receipt retains the action, provider
event identifier, provider status, and reviewed scope hash, but no access token or
raw provider response.

If an adapter invocation fails after transmission, the action becomes
the ``execution unknown; reconciliation required'' state; it is not automatically
repeated.
This prevents an uncertain response from creating duplicate events or applying a
modification twice.

\subsection{Google Calendar Integration}

The implementation contains provider adapters for event insertion, exact
rescheduling, and deletion, plus a private free/busy contract.  Real execution is
disabled until the active persona completes PKCE authorization and grants the
reviewed Calendar scopes.  The account binding belongs to that persona and cannot
be selected by shared-TV speech.

\section{Pilot Guardian: First Safe Urgent-Help Proof}
\label{sec:pilot-guardian}

Pilot Guardian is a non-medical urgent-help system.  Its first proof connects an
explicit spoken or phone request to a large confirmation surface, one or more
separately authorized trusted contacts, optional consented location, and a live
response record.  It does not diagnose a fall and does not represent itself as an
ambulance or emergency-service dispatcher.

\subsection{Permission Separation}

An ordinary contact permission does not imply emergency interruption.  The owner
must grant a separate Guardian permission naming the contact, route, interruption
policy, and whether location may be shared.  Location requires its own explicit
permission receipt.  Entry, lock, light, wearable, and room-sensor permissions are
not inferred from contact status.

\subsection{Incident Flow}

The implemented state machine is
\begin{enumerate}
  \item an explicit help phrase or private-phone control creates an incident;
  \item the TV and phone show a large, cancellable confirmation state;
  \item the protected identity confirms or marks the event as a false alarm;
  \item confirmation selects only pre-authorized Guardian contacts;
  \item an authorized adapter attempts each alert and must return a provider alert ID;
  \item failed routes become a visible connectivity-fallback state;
  \item a trusted contact may respond through the incident's unguessable response token;
  \item the response is shown as live human contact evidence.
\end{enumerate}

At every state the structured record contains
the ambulance-dispatched and emergency-service-contacted flags set to false,
unless a future tested official
adapter proves otherwise.  No language-model output can change these fields.

\subsection{Sensors and False Positives}

Authorized wearable or room-sensor observations may be recorded as
\texttt{possible\_fall}, \texttt{wearable\_sos}, or \texttt{inactivity} signals.
They include confidence and provenance but explicitly set
the medical-diagnosis and automatic-dispatch flags set to false.  A
signal prompts the human confirmation flow.  The protected identity can cancel an
unconfirmed incident as a false alarm.

\subsection{External Safety Work Still Required}

Guardian is not field-complete.  Production deployment requires regional legal,
clinical, accessibility, safeguarding, privacy, telecoms, and emergency-service
review; resilient push/SMS/voice providers; power and Internet failure testing;
continuous delivery-path monitoring; escalation timing studies; false-positive
field trials; and explicit qualification in every supported country.  In Spain
and the European Union, an ordinary app must not claim connection to the 112
emergency system without an authorized, tested integration and verified receipt.

\section{Verification and Release Boundary}

Automated tests cover local-first drafting, exact recipient/content hashes,
guardian approval for child communication, provider receipt requirements,
received-message provenance, message-to-task conversion, conditional no-reply
follow-up, calendar conflict privacy, explicit conflict acceptance, free-slot
queries, create/reschedule/cancel receipts, PKCE state and scope validation, opaque
token-vault references, Guardian permissions, optional location, false-alarm
cancellation, sensor non-diagnosis, connectivity fallback, verified trusted-contact
alerts, two-way response records, and deterministic emergency voice routing.

The locally enforceable software is complete only at its adapter boundaries.
Real Gmail, Calendar, messaging, call, notification, and emergency delivery remain
disabled until the owner or product operator supplies registered provider clients,
completes OAuth or equivalent authorization, and obtains the required production
and regional approvals.

\subsection{Official Technical References}

The implementation follows Google's installed-application OAuth and PKCE guidance,
Gmail \texttt{users.messages.send}, Calendar \texttt{events.insert}, exact event
modification APIs, and Calendar free/busy query.  These provider interfaces and
policies must be revalidated before production release because authorization
requirements and platform policies can change.
